\section{Module Names}\label{sec:Modules}
A module\index{module} is defined in the file \verb#module-info.java#\index{module-info.java} that is located in the root of the source tree for that module\index{module}; see chapter \tqfullvref{sec:ModuleDefinition} and \autocite{ORACLE_DOC_LANGUAGE_SPECIFICATION:ModuleDeclarations} for the details.

Each module\index{module} has a main package\index{package!main package}, and the full name of this package\index{package} should be also the name of the module\index{module}. If the module\index{module} will export any packages\index{package} at all\footnote{A module\index{module} for a program is not required to export anything.}, it has to export at least this main package\index{package!main package}.

So the names for modules\index{module} will follow basically the same rules as those for packages\index{package}, as described in chapter \tqref{sec:Packages}.

%==============================================================================
% End of file!!
%==============================================================================
